\section{Introduction}
\label{sec:intro}

The proliferation of web-based food delivery platforms has revolutionized how customers discover, order, and receive meals from restaurants. Driven by changing consumer preferences, the rapid growth of e-commerce and mobile technologies has expanded access to diverse cuisines, streamlined restaurant operations, and provided new opportunities for digital innovation~\cite{b7,b8,b9}. Food delivery services now form a critical segment of the global e-commerce market, where reliability, usability, and efficiency are key determinants of user satisfaction and business success.

Major platforms such as Zomato, Swiggy, DoorDash, and UberEats have popularized integrated online ordering through advanced recommendation systems, real-time logistics, and large-scale marketplace management~\cite{b10}. Despite their popularity, these platforms face several fundamental challenges that limit system reliability and scalability. Common issues include inaccurate inventory synchronization, delayed cart updates during high-traffic periods, incomplete order processing, and inconsistent user experience across devices~\cite{b11}. Moreover, traditional monolithic architectures in food delivery services struggle to balance transaction reliability with flexibility for multi-restaurant environments, resulting in periodic stock mismatches and elevated rates of cart abandonment~\cite{b1}.

The increasing demand for modular, scalable systems is matched by a need for robust session management and data consistency~\cite{b12}. User expectations for seamless transitions between device sessions, real-time cart updates, and responsive browsing necessitate architectural solutions that minimize latency while maintaining data integrity. Food delivery platforms must address key technical requirements: 1) multi-user inventory accuracy, 2) consistent cart state across sessions and devices, 3) scalable order management workflows, and 4) secure authentication for both anonymous and registered users~\cite{b2,b13}.

\textbf{PORTKEY} is developed in response to these needs, embodying a session-aware web application that prioritizes inventory management, cart persistence, and modular design. Implemented using Flask and SQLAlchemy ORM with SQLite for rapid prototyping and reliable transaction support, PORTKEY demonstrates that well-architected web applications can efficiently address the persistent challenges facing food delivery e-commerce systems.

The system adopts a hybrid session management strategy, supporting both anonymous and logged-in users, and ensures inventory accuracy through atomic cart operations~\cite{b4}. Real-time inventory tracking and seamless cart state synchronization are integrated as core platform features, ensuring the prevention of overselling and enhancing overall user satisfaction~\cite{b2}. The architecture embraces modular implementation principles, clear separation of concerns, and RESTful API design to guarantee maintainability, scalability, and extensibility for future research and production use~\cite{b5,b14}.

This paper contributes to the literature by presenting a unified and production-viable approach to web-based food ordering application architecture, focusing on practical solutions for session management, inventory coordination, and database normalization~\cite{b15}. While advanced features such as machine learning recommendations, multi-language support, and payment gateway integration are not included in the current implementation, the foundational patterns outlined in PORTKEY provide robust and extensible groundwork for subsequent development and research.

% Additional references for the introduction
\bibitem{b7} F. Oliveira et al., "A Thematic Review on Using Food Delivery Services during the Pandemic: Insights for the Post-COVID-19 Era," Int. J. Environ. Res. Public Health, vol. 19, no. 22, 2022.
\bibitem{b8} Y. Lee and D. Kim, “Enhancing Food Delivery with Chatbot-Assisted Ordering,” IEEE Internet Computing, vol. 28, no. 3, pp. 22–31, 2023.
\bibitem{b9} H. Patel, G. Mehta, "Contextual Recommendation Systems for Food Services," Expert Systems with Applications, vol. 233, 2024.
\bibitem{b10} S. Chen, P. Wang, “Optimizing Customer Repurchase Intention through Cognitive and Affective Experience: An Insight of Food Delivery Applications,” Sustainability, vol. 14, no. 19, 2022.
\bibitem{b11} T. Sharma, "A Review of the Usable Food Delivery Apps," Int. J. Eng. Res. Technol., vol. 8, no. 12, pp. 328–335, 2019.
\bibitem{b12} HCL Commerce, "Session management," 2022.
\bibitem{b13} Feedonomics, "Inventory Sync for Multichannel Ecommerce: How It Works," 2025.
\bibitem{b14} Pratap Sharma, “Architecture and Design Principle for Online Food Delivery,” 2022.
\bibitem{b15} GeeksforGeeks, “Database Design for Online Restaurant Reservation and Food Delivery,” 2024.
